\begin{figure}[p]
\small
\[
\begin{array}{rcl}
\ell&::=&\lo\mid\hi\qquad(\lo\sqsubseteq\hi)\\
e&::=&s\mid x\mid e\mathbin{+\!+}e\mid\mathsf{escape}(e)\mid\mathsf{trim}(e)\\
g&::=&x=s\mid\mathsf{prefix}(x,s)\\
\pi&::=&\bot\mid(T,N)\\
p&::=&\mathsf{skip}\mid p;p\\
 &&\mid\mathsf{call}_{i}(x,m,k,e,\pi)\mid\mathsf{tool}(x,f,e)\\
 &&\mid\mathsf{dummy}(m,T,N)\mid\mathsf{dummyTool}(f)\\
 &&\mid\mathsf{if}\ g\ \mathsf{then}\ p\ \mathsf{else}\ p\\
 &&\mid\mathsf{loop}(n,g,p,b),\qquad b\in\{\mathsf{false},\mathsf{true}\}.
\end{array}
\]
Strings are $s$; variables and randomness sites are natural numbers $x,i$.
The model contract is $m$, the useful output cap is $k$, and padding annotation
$(T,N)$ fixes logical input and output counts. The loop runs at most $n$ times,
stops when $g$ holds, and emits cover iterations when $b=\mathsf{true}$.
This grammar lists every constructor of Lean's \code{Expr}, \code{Guard} and
\code{Stmt}; $\pi=\bot$ denotes an unpadded call.

The corresponding surface fragments are
\begin{lstlisting}[basicstyle=\ttfamily\footnotesize]
input x: text secret max_tokens c;
prompt P(x) = user "template {x}";
let y = call P(x) using m max_tokens k;
let y = tool f(x);
if x == "value" { ... } else { ... }
repeat n until x == "done" { ... }
\end{lstlisting}
The trusted frontend expands templates into expressions, allocates slots and
sites, and joins labels into a fixed context $\Gamma$. Literal model/template
framing is part of the request expression rather than an implicit token charge.
\caption{Complete core grammar and representative surface constructs. Surface
parsing and lowering are tested implementations outside the Lean proof.}
\label{fig:grammar}
\end{figure}
